\section{Involución entre las coordenadas de residuo y cociente}
\label{sec:dualidad-visible-memoria}

La división euclídea usada en \(\mathcal A_9^\#\) escribe, para un entero en
la carta nonádica,

\[
n=9Q+r,
\qquad r\in\{1,\ldots,9\}.
\]

La coordenada \(r\) es el representante observable y \(Q\) registra el
número de cruces de frontera. Para formular una involución que intercambie
ambas coordenadas es necesario pasar a una extensión estable bajo el
intercambio. Denotemos por \(\bar r\in\{0,\ldots,8\}\) la clase residual,
con \(\bar r=0\) cuando el representante es \(r=9\), y consideremos

\[
\mathscr D_{\mathrm{dual}}=\mathbb Z^2\times\mathbb R_{>0}.
\]

Los dos enteros de \(\mathscr D_{\mathrm{dual}}\) no se restringen a la sección canónica de
representantes; esta ampliación es necesaria porque un cociente \(Q\) puede
quedar fuera del intervalo \(0,\ldots,8\) al ocupar la primera coordenada.

\begin{definition}[Forma de energía de radio]
Para \(R_0>0\) y \((m,w)\in\mathbb Z^2\), definimos
\[
E_{R_0}(m,w)=\frac{m^2}{R_0^2}+w^2R_0^2.
\]
La transformación de dualidad es
\[
\mathsf T(m,w,R_0)=(w,m,R_0^{-1}).
\]
\end{definition}

\begin{theorem}[Dualidad residuo--cociente]
\label{thm:dualidad-visible-memoria}
La aplicación \(\mathsf T\) es una involución de \(\mathscr D_{\mathrm{dual}}\) y preserva
la forma de energía en el sentido
\[
E_{R_0}(m,w)=E_{R_0^{-1}}(w,m).
\]
\end{theorem}

\begin{proof}
La segunda aplicación de \(\mathsf T\) intercambia de nuevo \(m,w\) e
invierte por segunda vez \(R_0\), por lo que
\[
\mathsf T^2(m,w,R_0)=(m,w,R_0).
\]
Además,
\[
E_{R_0^{-1}}(w,m)
=\frac{w^2}{R_0^{-2}}+m^2R_0^{-2}
=w^2R_0^2+\frac{m^2}{R_0^2}
=E_{R_0}(m,w).\qedhere
\]
\end{proof}

La inserción de un estado nonádico en el dominio ampliado es

\[
(r,Q)\longmapsto(\bar r,Q,R_0).
\]

El transformado \((Q,\bar r,R_0^{-1})\) no tiene por qué pertenecer a la
sección \(0\leq\bar r\leq8\). Por esta razón, \(\mathsf T\) es una
involución exacta del dominio ampliado, no un endomorfismo del atlas residual
de 81 posiciones. El punto fijo de la transformación del parámetro es
\(R_0=1\); allí el intercambio de las dos coordenadas preserva la misma forma
cuadrática.

La identidad admite una realización operatorial. Sea
\(\mathcal H_{\rm lat}=\ell^2(\mathbb Z^2)\), con base ortonormal
\(\{|m,w\rangle\}\). Definamos el operador diagonal

\[
H(R_0)|m,w\rangle=E_{R_0}(m,w)|m,w\rangle
\]

en su dominio natural, y el operador unitario

\[
U_T|m,w\rangle=|w,m\rangle.
\]

\begin{corollary}[Equivalencia unitaria de radios recíprocos]
\label{cor:dualidad-unitaria}
Para todo \(R_0>0\),
\[
U_TH(R_0)U_T^{-1}=H(R_0^{-1}).
\]
\end{corollary}

\begin{proof}
Sobre un vector de la base,
\[
U_TH(R_0)U_T^{-1}|m,w\rangle
=E_{R_0}(w,m)|m,w\rangle.
\]
El \cref{thm:dualidad-visible-memoria} da
\(E_{R_0}(w,m)=E_{R_0^{-1}}(m,w)\), que es la acción de
\(H(R_0^{-1})\).
\end{proof}

Esta conjugación tiene la forma matemática de la dualidad de radio en una
compactificación circular, en unidades en las que la escala de cuerda se ha
normalizado a uno \cite{Giveon1994}. La identificación física adicional
asigna \(m\) al número de momento, \(w\) al número de enrollamiento y \(R_0\)
al radio de compactificación. El teorema no introduce osciladores, campos,
supersimetría ni una tensión de cuerda; demuestra la involución reticular y la
equivalencia unitaria de los operadores diagonales declarados.
